\chapter{Project foundation}

Creating a good foundation

\section{Setting up the project}

Not as easy as pressing File > New Project

As this is going to be a project that multiple people work on we need a way to synchronies our work. For this

\subsection{Build systems}

To give the project a good foundation and to even start development I had to chose a build system. As C/C++ doesn't have a standard build infrastructure, there are many different ones that have their own pros, and cons.

\subsubsection{Msbuild}
The obvious choice
MsBuild is the build system Microsoft uses in VisualStudio and what most of the games companies use.
It is an obvious choice as it has great integration with VS and the Windows ecosystem.
But it also locks the project to be windows only, as it is nearly impossible to build a MSBuild C++ project on Linux or MacOS.
This means that we can't use it if we ever want to make the engine multi-platform.
MSBuild also locks us in to Visual Studio as the code editor, and not every potential contributor likes it, so this is also a problem.
As an open source project we want to allow contributors to use whatever IDE they want to minimizes friction, and make onboarding as easy as possible.

\subsubsection{Bazel}
The new kid
Bazel is a relatively less used build system that Google developed for their own monorepo.
It is optimized for large projects that have complex interdependencies.
For us the two major features that would be useful are the build cache and multiple language support.
The Bazel build cache is capable of uploading to a server the results of individual compilation steps and then allow all other developer to download the results.
This is accomplished by making the build steps strictly deterministic and caching the result.
The second important thing Bazel can offer is multiple language support,
this means that it would be possible to truly have a single repository for all that is connected to teh project.
For example the repo could contain the website, thesis latex and engine code.

Problems with bazel

Lack of language servers
Bad integration

\subsubsection{CMake}
The common one

CMake is the most common build system among cross-platform projects.
It is not a single ready to use package tho.
Users can generate multiple types of configurations from the same CMake files.
This means that on Windows one could generate Visual Studio project files, on Linux Make files and so on.

\subsubsection{The choice}
I have experimented and tried the project setup with all 3 of the options that came up.
MsBuild was quickly discarded as none of us wanted to use Visual Studio.
It turned out that our preferred editor is CLion from Intelij.
This meant that both Bazel and Cmake could work.

First I tried to use Bazel but that turned out to be lacking many comforts that a modern coding environment is expected to have (See Bazel>Problems).
Because of this I even had a call with the current product manager of Bazel where I described most of the issues I had, I hope to try it again and be more successful in a year or so.

Finally, we settled on CMake as it is a known usable system and has excellent integration in CLion.
~~This is also a known tool in the Linux community and should allow us to build on any platform.~~

One of the important things CMake allows us to do is minimize the amount of dependencies the developers need to manually download.
This means that all the libraries that we use are downloaded by CMake based on versioned URLs